%-------------------------------------------------------------------- Comente: 1er vs 3er grado
\needspace{8\baselineskip}
\tema{Comente: discriminación de primer vs tercer grado}
\question \textbf{Comente.} ¿Cuál es la diferencia entre discriminación de primer y tercer grado?
\needspace{5\baselineskip}
\begin{solution}
En la discriminación de primer grado la firma puede fijar el precio que maximice su excedente, dejando nada del excedente al consumidor (cobra a cada consumidor su precio de reserva). En la discriminación de tercer grado la firma explota las características observables del comprador para cobrar precios diferenciados según el grupo.

\end{solution}

%-------------------------------------------------------------------- Comente: mercados
\needspace{8\baselineskip}
\tema{Comente: mercados con discriminación de primer y tercer grado}
\question \textbf{Comente.} Caracterice los mercados en los que suele haber discriminación de primer grado y de tercer grado.
\needspace{5\baselineskip}
\begin{solution}
Los mercados en donde suele haber discriminación de primer grado son:
\begin{itemize}
  \item Mercados donde el número de compradores (\emph{customers}) es relativamente pequeño y el vendedor posee considerable información sobre los compradores.
  \item Industrias como el concreto fresco, aviones de pasajeros, software especializado para empresas, etc.
  \item A pesar de que existe un precio de lista, cada cliente recibe un descuento que se negocia: el precio final depende de la disposición a pagar del cliente y de su poder de negociación.
\end{itemize}
Por otro lado, los mercados donde suele haber discriminación de tercer grado son mercados más grandes, donde los costos de información son mayores; por ejemplo, los pasajes del transporte público.

\end{solution}

%-------------------------------------------------------------------- Comente: captar todo el excedente
\needspace{8\baselineskip}
\tema{Comente: captar todo el excedente del consumidor}
\question \textbf{Comente.} Es imposible que una firma monopólica logre captar todo el excedente del consumidor, puesto que los precios y la estrategia que impongan dependerán también de factores como la elasticidad de la demanda.
\needspace{5\baselineskip}
\begin{solution}
Es verdad que en la práctica es difícil que se logre captar todo el excedente, pero si el monopolio tiene información perfecta sobre la disposición a pagar de los consumidores podría imponer una estrategia de discriminación de precios de primer grado, donde cada individuo paga el máximo que está dispuesto a pagar por el producto en cuestión, generando que el excedente sea nulo para los consumidores y que la firma se lleve todo el excedente del mercado.

\end{solution}

%-------------------------------------------------------------------- Comente: discriminar o no
\needspace{8\baselineskip}
\tema{Comente: ¿discriminar o no discriminar?}
\question \textbf{Comente.} ¿Discriminar o no discriminar? Discuta los principios básicos para responder esta pregunta.
\needspace{5\baselineskip}
\begin{solution}
Cualquier baja en la producción suele llevar a que el bienestar total baje. Una discriminación que aumenta la demanda del producto o servicio suele aumentar el bienestar (no hay consumidor más triste que el que no consume). Los principios básicos son:
\begin{itemize}
  \item Si la producción total cae con la discriminación, el bienestar total disminuye.
  \item Si un monopolista que no discrimina cierra un mercado, es mejor discriminar.
\end{itemize}

\end{solution}

%-------------------------------------------------------------------- Museo
\needspace{8\baselineskip}
\tema{Discriminación de tercer grado: el museo}
\question En un pueblo del sur de Chile hay un único museo que recibe a visitantes nacionales y extranjeros, quienes presentan una mayor valoración del museo. El costo marginal de producción del museo es 1 por cada visitante y no hay costos fijos. Las demandas de extranjeros y nacionales son
\begin{align*}
  Q_E &= 10 - P_E \\
  Q_N &= 8 - P_N
\end{align*}
\begin{parts}
\part Suponga que el museo monopólico fija un precio uniforme. ¿Qué precio cobra y qué beneficio obtiene? ¿Qué condición debe cumplirse para que se venda a ambos grupos?
\needspace{5\baselineskip}
\begin{solution}
La firma maximiza beneficios sobre la demanda agregada, $Q_T = 18 - 2P$:
\[
\begin{gathered}
  \max_P\ \pi = (P - c)(18 - 2P) \\
  \;\Rightarrow\; 18 - 4P + 2c = 0 \\
  \;\Rightarrow\; P^U = 5, \\
  \qquad \pi^U = (5-1)(18 - 2\cdot 5) = 32 .
\end{gathered}
\]
Para que se venda a ambos grupos debe cumplirse la restricción de participación del grupo de menor valoración: los nacionales no demandarían ninguna entrada a precios mayores a 8, por lo que debe cumplirse $P \leq 8$ (la demanda agregada $18-2P$ solo es válida en ese tramo). Como $P^U = 5 < 8$, se vende a ambos.

\end{solution}

\part Suponga ahora que el monopolio puede discriminar en tercer grado. Calcule los precios de cada mercado y los beneficios del museo.
\needspace{5\baselineskip}
\begin{solution}
Para los extranjeros el museo resuelve
\[
\begin{gathered}
  \max_{P_E}\ (P_E - c)(10 - P_E) \\
  \;\Rightarrow\; 10 - 2P_E + c = 0 \\
  \;\Rightarrow\; P_E = \frac{10 + c}{2} = 5{,}5, \qquad \pi_E = (5{,}5-1)(10-5{,}5) = 20{,}25 .
\end{gathered}
\]
Y de la misma manera, para los nacionales,
\[
  \max_{P_N}\ (P_N - c)(8 - P_N)
  \;\Rightarrow\; P_N = \frac{8 + c}{2} = 4{,}5, \qquad \pi_N = (4{,}5-1)(8-4{,}5) = 12{,}25 .
\]
Los beneficios totales son $\pi^D = \pi_E + \pi_N = 32{,}5$.

\end{solution}

\part Un compañero le menciona que el museo debería enfocarse solo en el grupo de alta valoración. Demuéstrele que al museo no le conviene enfocarse solo en una parte del mercado, cuando tiene la opción de discriminar precios.
\needspace{5\baselineskip}
\begin{solution}
Si el museo se enfocara solo en los extranjeros, fijaría un precio en que solo ellos demanden una cantidad positiva ($P = 8$). Los beneficios serían entonces
\[
  \pi_E = (8-1)(10-8) = 14,
\]
lo cual es bastante menor que al no discriminar ($32$) o al discriminar en tercer grado ($32{,}5$).

\end{solution}

\part Comparando los beneficios del monopolio con la estrategia de precio uniforme y de discriminación, ¿cuál le conviene más? Además, calcule los excedentes de los consumidores extranjeros y nacionales. ¿Cuál grupo se beneficia de la discriminación de tercer grado y cuál se perjudica?
\needspace{5\baselineskip}
\begin{solution}
Como $\pi^D = 32{,}5 > \pi^U = 32$, al museo le conviene discriminar en tercer grado (aunque la ganancia es marginal).

Los excedentes de los consumidores con precio uniforme (U) y con precio diferenciado (D) son:
\begin{align*}
  EC_E^U &= \frac{5(10-5)}{2} = 12{,}5, & EC_E^D &= \frac{4{,}5(10-5{,}5)}{2} = 10{,}125, \\
  EC_N^U &= \frac{3(8-5)}{2} = 4{,}5, & EC_N^D &= \frac{3{,}5(8-4{,}5)}{2} = 6{,}125 .
\end{align*}
El grupo de mayor valoración (extranjeros) sale perjudicado, mientras que el de menor valoración (nacionales) está mejor siendo discriminado. En cuanto al excedente total, los consumidores están peor con discriminación ($17 \to 16{,}25$).

\end{solution}

\part ¿Qué pasa con el beneficio social? Calcule también el producto total ofrecido en cada estrategia y comente su relación con los cambios en bienestar o ineficiencias que puedan estar ocurriendo. ¿Qué otros factores explican los cambios en bienestar al pasar de una estrategia a otra?
\needspace{5\baselineskip}
\begin{solution}
El beneficio social para cada estrategia corresponde a
\[
  ES^U = \pi^U + EC^U = 32 + 17 = 49, \qquad ES^D = \pi^D + EC^D = 32{,}5 + 16{,}25 = 48{,}75 .
\]
Se concluye que en este caso el bienestar social no aumenta con la discriminación de precios. La producción total de cada estrategia es
\[
  Q^U = 18 - 2\cdot 5 = 8, \qquad Q^D = Q_E^D + Q_N^D = (10 - 5{,}5) + (8 - 4{,}5) = 4{,}5 + 3{,}5 = 8 .
\]
Como el producto total no cambia, no hay un efecto de expansión de la producción que compense; el bienestar cae por una \emph{mala asignación} del bien entre los consumidores: con precios distintos hay consumidores con mayor valoración que consumen menos y consumidores con menor valoración que consumen más (se pierde eficiencia en la asignación).

\end{solution}
\end{parts}

%-------------------------------------------------------------------- Comente: condiciones discriminación
\needspace{8\baselineskip}
\tema{Comente: condiciones para discriminar precios}
\question \textbf{Comente.} La discriminación de precios solo es posible si es que se cumplen las siguientes condiciones: (i) es posible establecer que existen grupos con diferentes demandas o disposiciones a pagar en el mercado y (ii) es posible identificar explícitamente a los individuos que pertenecen a cada grupo.
\needspace{5\baselineskip}
\begin{solution}
Falso. La discriminación de precios de segundo grado no requiere que se cumpla la condición (ii). Se pueden diseñar ofertas o planes tales que cada consumidor se autoseleccione en el servicio que va dirigido a su propio grupo.

\end{solution}

%-------------------------------------------------------------------- Tarifas en dos partes
\needspace{8\baselineskip}
\tema{Tarifas en dos partes (dos tipos de clientes)}
\question Una empresa ofrece dos planes de servicio para sus clientes, los que consisten en un monto fijo $F$ y un precio $P$ por unidad. Esta empresa se enfrenta a dos mercados distintos con las siguientes funciones de demanda:
\begin{align*}
  Q_1 &= 100 - P \\
  Q_2 &= 80 - P
\end{align*}
Supondremos que ofrecer el servicio a cualquier mercado tiene un costo de cero. Además, la firma tiene la opción de fijar una tarifa fija $F$ además de una tarifa por unidad $P$.
\begin{parts}
\part Si usted puede separar a los clientes de cada mercado, ¿cuál sería la tarifa $F$ y cuál la tarifa por unidad $P$ para cada tipo de cliente? ¿Cuánto sería el beneficio de la empresa?
\needspace{5\baselineskip}
\begin{solution}
Si puedo separar a los clientes de cada mercado significa que sé quién es quién (y de esa forma puedo evitar el arbitraje). Lo más lógico es extraer a cada individuo su máxima disposición a pagar, que es igual al excedente máximo de cada individuo. El $F$ óptimo busca apropiarse del excedente del consumidor que queda al cobrar una tarifa por unidad $P$:
\[
  F_1 = \frac{(100 - P_1)^2}{2}, \qquad \pi_1 = F_1 + P_1(100 - P_1) = \frac{(100-P_1)^2}{2} + (100 - P_1)P_1 .
\]
Esta expresión se maximiza con $P_1 = 0$ (cobrar al costo marginal, para maximizar el excedente total que luego se extrae), por lo tanto $F_1 = \frac{100^2}{2} = 5000$. Con la misma lógica, $P_2 = 0$ y
\[
  F_2 = \frac{80^2}{2} = 3200 .
\]
Los beneficios totales de la empresa serán $\pi = 5000 + 3200 = 8200$.

\end{solution}

\part Suponga que no puede separar a los clientes y la tarifa fija $F$ está restringida a cero. ¿Cuál sería la tarifa $P$ que cobraría? ¿Cuánto sería el beneficio de la empresa?
\needspace{5\baselineskip}
\begin{solution}
Si no se puede separar a los clientes, se cobrará un precio uniforme en ambos mercados. Como $F=0$, solo se debe decidir $P$. Sumando las demandas horizontalmente, la demanda total es
\[
  D(P) =
  \begin{cases}
    100 - P & \text{si } P > 80 \\
    180 - 2P & \text{si } P \leq 80 .
  \end{cases}
\]
El tramo donde maximizará la firma es el segundo (si se maximiza en el primero se obtiene $P = 50$, pero la condición del tramo es $P>80$):
\[
  \max_P\ \pi = (180 - 2P)P \;\Rightarrow\; P = 45, \quad Q = 90, \quad \pi = 4050 .
\]
Los beneficios son menores que en la parte (a).

\end{solution}

\part Suponga ahora que cobra una tarifa en dos partes, es decir, $F$ y $P$, pero sin diferenciar al tipo de cliente, sino que para todos lo mismo. ¿Cuál sería la tarifa fija $F$ y cuál la tarifa por unidad $P$? ¿Cuánto sería el beneficio de la empresa?
\needspace{5\baselineskip}
\begin{solution}
En este caso no podemos cobrar una tarifa fija de $5000$ como en el primer caso, pues los clientes de demanda baja no consumirían. Por lo tanto, cobramos la tarifa fija más alta posible que sigue siendo aceptada por ambos, que es precisamente el excedente de los consumidores de demanda baja:
\[
  F = \frac{(80 - P)^2}{2}.
\]
La empresa maximiza
\[
\begin{gathered}
  \max_P\ \pi = 2F + P(100 - P) + P(80 - P) = (80-P)^2 + P(180 - 2P) \\
  \;\Rightarrow\; \\
  \frac{\partial\pi}{\partial P} = -2(80-P) + 180 - 4P = 20 - 2P = 0 \\
  \;\Rightarrow\; P = 10 .
\end{gathered}
\]
La tarifa fija correspondiente es $F = \frac{(80-10)^2}{2} = 2450$, y los beneficios son $\pi = 2\cdot 2450 + 10\cdot(90 + 70) = 6500$.

\end{solution}

\part ¿De qué forma la tarifa en dos partes le permite discriminar? ¿Qué restricciones enfrenta al maximizar beneficios esta empresa?
\needspace{5\baselineskip}
\begin{solution}
El beneficio con tarifa en dos partes ($6500$) es mayor que cobrar solo un precio unitario ($4050$). Esto se produce porque a través de la tarifa en dos partes se puede inducir a los distintos tipos de individuos presentes en el mercado (demanda alta vs.\ demanda baja) a autoseleccionarse de acuerdo a los menús de precios y tarifas ofrecidos. Esta discriminación será posible si los individuos presentes en el mercado tienen demandas distintas. Con respecto a las restricciones que enfrenta la firma a la hora de maximizar, existen dos tipos:
\begin{enumerate}
  \item[(a)] \emph{Restricción de participación:} la tarifa y el precio deben ser lo suficientemente atractivos para que los individuos a los cuales se quiere vender, compren.
  \item[(b)] \emph{Restricción de compatibilidad de incentivos:} puede que los individuos de demanda alta tengan el incentivo a comprar el paquete de menor costo. Por lo tanto, se debe asegurar que los clientes no se hagan pasar unos por otros.
\end{enumerate}

\end{solution}
\end{parts}

%-------------------------------------------------------------------- Comente: 2º grado tarifa en dos partes
\needspace{8\baselineskip}
\tema{Comente: discriminación de segundo grado}
\question \textbf{Comente.} En una discriminación de segundo grado el monopolista no observa ninguna característica del consumidor que le permita aplicar una tarifa en dos partes similar a la discriminación perfecta.
\needspace{5\baselineskip}
\begin{solution}
Verdadero. Bajo discriminación de segundo grado el monopolista no observa directamente alguna característica que le permita separar a los consumidores, pero puede establecer una tarifa en dos partes con un cargo variable superior al costo marginal y un cargo fijo equivalente al excedente que obtendría el consumidor con menor disponibilidad a pagar. Esta tarifa difiere de la aplicada bajo discriminación perfecta en que allí el cargo variable es igual al costo marginal y el cargo fijo puede ser diferenciado.

\end{solution}

%-------------------------------------------------------------------- Comente: esquemas de tarifación
\needspace{8\baselineskip}
\tema{Comente: tipos de discriminación (cine, buffet, transporte)}
\question \textbf{Comente.} Es interesante observar que los esquemas de tarificación difieren en distintas actividades. Describa el tipo de discriminación de precios y por qué se usa (o no se usa) en los siguientes casos:
\begin{parts}
\part Un cine durante un día normal.
\needspace{5\baselineskip}
\begin{solution}
Al haber descuentos para estudiantes y tercera edad, esto constituye una discriminación de tercer grado. Con las salas \emph{premium} de algunos cines se ofrecen paquetes distintos, lo que es discriminación de segundo grado para que los clientes se autoseleccionen según su disposición a pagar.

\end{solution}

\part Restaurantes con buffet (se puede comer cuanto se desea).
\needspace{5\baselineskip}
\begin{solution}
Discriminación de segundo grado. Los clientes se autoseleccionan; en general se cobra una tarifa fija porque no se puede controlar el consumo.

\end{solution}

\part Transporte público con pasajes especiales para estudiantes.
\needspace{5\baselineskip}
\begin{solution}
Discriminación de tercer grado, pues los pases escolares de los estudiantes son observables.

\end{solution}
\end{parts}

%-------------------------------------------------------------------- Comente: plan de telefonía
\needspace{8\baselineskip}
\tema{Comente: planes de telefonía móvil}
\question \textbf{Comente.} Suponga que usted es una/un feliz ayudante de la FEN-UChile que se dispone a recibir su pago de ayudantía en los próximos días. Usted ha decidido destinar este dinero a renovar su celular, el que actualmente tiene un plan de 200 minutos e internet ilimitado. Con esta idea se dirige a la compañía de telefonía móvil y cotiza un nuevo equipo. La vendedora le hace el siguiente comentario: «El nuevo equipo cuesta \$50.000. Sin embargo, su plan actual de \$25.990 está descontinuado. Dado esto, lo homologaremos al más cercano de los nuevos sin perjudicar su pago. Este corresponde al plan de \$23.990, que consta de 150 minutos y 1 Gb de internet. Adicionalmente, le ofrecemos el plan de \$29.990 con 300 minutos y 2,5 Gb de internet». ¿Qué estrategia está utilizando la compañía móvil? ¿De qué forma funciona?
\needspace{5\baselineskip}
\begin{solution}
La compañía móvil está utilizando una estrategia de discriminación de segundo grado. La empresa no puede exigir que el consumidor esté en cierto plan, pero ofrece dos planes distintos donde el cliente tiene que autoseleccionarse. En particular, lo que hace la firma es perjudicar la opción que corresponde a los consumidores de menor valoración, disminuyendo el número de minutos y ofreciendo una cantidad limitada de internet móvil. Dado lo anterior, el excedente que puede obtener el individuo de mayor valoración escogiendo este plan disminuye. Esto puede empujar a que este decida optar por el plan más costoso.

\end{solution}

%-------------------------------------------------------------------- Comente: 3er grado bienestar
\needspace{8\baselineskip}
\tema{Comente: discriminación de tercer grado y bienestar}
\question \textbf{Comente.} Un economista afirma: «La discriminación de precios de tercer grado resulta en que algunos consumidores pagan un precio mayor que lo que hubieran pagado si el monopolista cobrara un precio uniforme. Por lo mismo, la discriminación de precios disminuye el bienestar y debería estar prohibida.» ¿Usted está de acuerdo con dicha afirmación? Fundamente su respuesta.
\needspace{5\baselineskip}
\begin{solution}
No estoy de acuerdo. La discriminación de precios puede aumentar o reducir el bienestar. Con discriminación de precios, los consumidores con mayor disposición a pagar (menor elasticidad de demanda) pagan precios mayores, pero los consumidores con menor disposición a pagar (mayor elasticidad de demanda) pagan precios menores. Al final del día, el bienestar puede aumentar o disminuir: depende de si la producción total aumenta (por ejemplo, si se abren mercados que no se servían) o no.

\end{solution}

%-------------------------------------------------------------------- Comente: aerolíneas
\needspace{8\baselineskip}
\tema{Comente: discriminación de precios en aerolíneas}
\question \textbf{Comente.} Debería prohibirse la discriminación de precios por parte de las aerolíneas, ya que esto les permite abusar de su poder de mercado.
\needspace{5\baselineskip}
\begin{solution}
Falso. En el extremo, incluso un monopolio discriminador en primer grado produce la cantidad socialmente óptima, por lo que en la medida que mejor pueda discriminar, disminuirá la pérdida social de eficiencia. Además, en la industria aérea existen costos fijos importantes, pero los costos variables y marginales son muy pequeños (una vez que el vuelo se programó y se pagó, el costo de subir pasajeros es prácticamente cero). En este caso, una efectiva discriminación de precios hace que las aerolíneas puedan mejorar sus factores de ocupación, disminuyendo sus costos medios y por ende pudiendo cobrar menores tarifas.

\end{solution}

%-------------------------------------------------------------------- AeroStrong
\needspace{8\baselineskip}
\tema{Discriminación de precios en aerolíneas: AeroStrong y AeroWeak}
\question La compañía aérea AeroStrong tiene dos tipos de pasajeros según la motivación del viaje: turismo y negocios. Las demandas de cada grupo son
\begin{align*}
  Q_t &= 20 - 4P_t \\
  Q_n &= 22 - P_n
\end{align*}
Los costos de movilizar un pasajero son independientes del motivo del viaje y es \$2.
\begin{parts}
\part Si la aerolínea puede evitar el arbitraje entre viajeros, ¿cuánto le cobraría por el pasaje a cada tipo de pasajero? En términos prácticos, ¿cómo la aerolínea podría evitar el arbitraje?
\needspace{5\baselineskip}
\begin{solution}
Si la empresa puede evitar el arbitraje entre los viajeros, cobrará un precio que maximice sus beneficios en cada mercado. Para los turistas:
\[
  \max_{P_t}\ (P_t - 2)(20 - 4P_t) \;\Rightarrow\; 28 - 8P_t = 0 \;\Rightarrow\; P_t = 3{,}5 .
\]
Para los pasajeros de negocios:
\[
  \max_{P_n}\ (P_n - 2)(22 - P_n) \;\Rightarrow\; 24 - 2P_n = 0 \;\Rightarrow\; P_n = 12 .
\]
Por ejemplo, si la aerolínea vende pasajes empaquetados ida y vuelta, podría identificar el propósito del viaje según si el pasajero permanece el fin de semana en el destino o no.

\end{solution}

\part Si la aerolínea se ve impedida de discriminar entre pasajeros, ¿cuál sería el precio que cobraría por los pasajes?
\needspace{5\baselineskip}
\begin{solution}
El monopolio fija un precio uniforme sobre la demanda agregada:
\[
  Q =
  \begin{cases}
    42 - 5P & \text{si } P \leq 5 \\
    22 - P & \text{si } P > 5
  \end{cases}
\]
El tramo donde maximizará la firma es el segundo: si maximizara en el primero el precio sería $P = 5{,}2$, pero la condición del tramo es $P \leq 5$. Entonces resuelve $\max_P\ (P-2)(22-P)$, cuya CPO es $24 - 2P = 0$, es decir, $P = 12$ (el grupo de turismo no consume).

\end{solution}

\part La aerolínea de bajo costo AeroWeak ingresa a competir en el segmento de pasajeros de turismo solamente, puesto que los pasajeros de negocios no la consideran una alternativa para sus viajes. El costo por pasajero de AeroWeak es \$1. ¿Qué es preferible para la entrante: que AeroStrong discrimine precios, que fije un precio único, o es indiferente dado que AeroWeak no compite en el segmento de negocios? Fundamente su respuesta.
\needspace{5\baselineskip}
\begin{solution}
Si AeroStrong fija un precio único ($12$), solamente abastecerá al grupo de negocios, por lo que AeroWeak podrá actuar como un monopolio en el grupo de turistas:
\[
\begin{gathered}
  \max_{P_t}\ (P_t - 1)(20 - 4P_t) \;\Rightarrow\; 24 - 8P_t = 0 \;\Rightarrow\; P_t = 3, \quad Q_t = 8, \\ \pi = 2\cdot 8 = 16 .
\end{gathered}
\]
Si AeroStrong discrimina precios, ambas empresas competirán en el segmento de turistas. Suponiendo que los turistas son indiferentes entre ambas aerolíneas y eligen la más barata, compiten à la Bertrand, y como AeroWeak es más eficiente fija $P_t = 2 - \epsilon$, con beneficios $\approx (2-1)\cdot(20-8) = 12 < 16$.

Por lo tanto, AeroWeak prefiere que AeroStrong fije un precio único, ya que le permite fijar un precio monopólico en el mercado de turistas.

\end{solution}
\end{parts}

%-------------------------------------------------------------------- Tarifas Clavistel
\needspace{8\baselineskip}
\tema{Tarifas en dos partes: Clavistel (variante numérica)}
\question La empresa de telefonía móvil Clavistel ofrece dos planes de servicio para sus clientes, los que consisten en un monto fijo $F$ y un precio $P$ por unidad del bien. El servicio a cualquier mercado tiene un costo de cero. Además, esta empresa se enfrenta a dos mercados distintos con las siguientes funciones de demanda: $q_1 = 40 - p$ y $q_2 = 20 - p$.
\begin{parts}
\part Si la empresa puede separar a los clientes de cada mercado, ¿cuál sería la tarifa $F$ y cuál la tarifa por unidad $P$ para cada tipo de cliente? Obtenga los beneficios de la empresa con esas tarifas.
\needspace{5\baselineskip}
\begin{solution}
Si la firma puede separar a sus clientes, le extraerá su máxima disposición a pagar, igual al excedente máximo de cada individuo. Para el cliente 1: $\frac{40\cdot 40}{2} = 800$. Con $F_1 = \frac{(40 - p_1)^2}{2}$, el beneficio es $F_1 + (40 - p_1)p_1$, que se maximiza en $p_1 = 0$ y $F_1 = 800$. Para el cliente 2, $p_2 = 0$ y $F_2 = \frac{20^2}{2} = 200$. Los beneficios totales son $800 + 200 = 1000$.

\end{solution}

\part Suponga ahora que la empresa no puede separar a los clientes y la tarifa fija $F$ está restringida a cero. ¿Cuál sería la tarifa $P$ que cobraría? ¿Cuánto sería el beneficio de la empresa?
\needspace{5\baselineskip}
\begin{solution}
Se cobra un precio uniforme en ambos mercados. La demanda total es $D(p) = 40 - p$ si $p > 20$ y $D(p) = 60 - 2p$ si $p \leq 20$. El tramo donde maximizará es el segundo (en el primero se tendría $p = 20$, justo en el borde, con beneficios menores): $\max_p\ (60 - 2p)p$, por lo tanto $P = 15$, $Q = 30$ y los beneficios son $450 < 1000$.

\end{solution}

\part Suponga ahora que cobra una tarifa en dos partes, es decir, $F$ y $P$, pero sin diferenciar al tipo de cliente. ¿Cuál sería la tarifa $F$ y $P$ que cobraría? ¿Cuánto sería el beneficio de la empresa?
\needspace{5\baselineskip}
\begin{solution}
No podemos cobrar una tarifa fija de 800, pues los clientes de demanda baja no consumirían; por lo tanto, cobramos la tarifa fija más alta posible, que es precisamente el excedente de los consumidores de demanda baja, $F = \frac{(20-p)^2}{2}$. La empresa maximiza
\[
\begin{gathered}
  \max_p\ \pi = 2\cdot\frac{(20-p)^2}{2} + (40 - p)p + (20 - p)p = (20-p)^2 + p(60 - 2p) \\
  \;\Rightarrow\; 20 - 2p = 0 \;\Rightarrow\; p = 10 .
\end{gathered}
\]
Entonces $F = \frac{10^2}{2} = 50$ y los beneficios son $\pi = 2\cdot 50 + 10\cdot(30 + 10) = 500$.

\end{solution}

\part ¿Qué restricciones enfrenta la empresa al maximizar beneficios discriminando con tarifa en dos partes?
\needspace{5\baselineskip}
\begin{solution}
La empresa enfrenta restricciones de dos tipos:
\begin{enumerate}
  \item[(a)] \emph{Restricción de participación:} debe hacer la tarifa y el precio lo suficientemente atractivos para que los individuos a los cuales se quiere vender, compren.
  \item[(b)] \emph{Restricción de autoselección (compatibilidad de incentivos):} no se debe permitir que los consumidores de un grupo usen los planes de otro.
\end{enumerate}

\end{solution}
\end{parts}

%-------------------------------------------------------------------- Tipos de discriminación
\needspace{8\baselineskip}
\tema{Los tres tipos de discriminación de precios}
\question Describa y explique los tres tipos de discriminación que puede hacer un monopolista según distintos casos.
\needspace{5\baselineskip}
\begin{solution}
\textbf{Segundo grado.} Discriminación por medio de la autoselección de los individuos según su comportamiento; por ejemplo, tarifas:
\begin{itemize}
  \item Tarifa lineal: $T(q) = pq$.
  \item Tarifa de dos partes: $T(q) = A + pq$, un cargo fijo más un cobro por unidad consumida.
  \item Tarifa no lineal: $T(q) = A_j + p_jq$ si $q_{j-1} < q < q_j$, con distintos cargos según tramos de consumo. Es común en bares, gimnasios, festivales de música, entre otros.
\end{itemize}
\textbf{Primer grado.} En el caso anterior el monopolista no conoce completamente la demanda por los individuos; si tuviera completo conocimiento sobre sus disposiciones a pagar, podría fijar tarifas que extrajeran todo su excedente. Asumiendo que no hay arbitraje, el monopolista podría tomar el precio de competencia ($p = c$) pero cobrar una suma fija según el individuo, que represente el excedente que obtendría:
\[
  T(q) =
  \begin{cases}
    A_i + cq & \text{si } q > 0 \\
    0 & \text{si no.}
  \end{cases}
\]
Se maximiza el bienestar social como en competencia perfecta, pero el monopolista se queda con todo el excedente. Esto suele estar presente en servicios más personalizados (\emph{customer markets}).

\textbf{Tercer grado.} A partir de información observable se pueden cobrar precios diferenciados. La información que tiene el monopolista es respecto a la demanda por grupos, pero no por individuo. El monopolista enfrenta un problema similar al del monopolista multiproducto:
\[
  \max_{p_1,\dots,p_n}\ \sum_{i=1}^{n}p_i q_i(p_i) - c\sum_{i=1}^{n}q_i(p_i).
\]
El ingreso marginal debe ser igual en cada grupo (con dos grupos, $i = 1,2$): $IMg_1 = c = IMg_2$. Desde el punto de vista de las elasticidades, $IMg_i = p_i\bigl(1 - \frac{1}{|\varepsilon_i|}\bigr) = c$, de modo que
\[
  p_i = \frac{c}{1 - 1/|\varepsilon_i|}.
\]
Mientras mayor (menor) sea la elasticidad en determinado grupo, se le cobra menos (más).

\end{solution}

%-------------------------------------------------------------------- Principios bienestar
\needspace{8\baselineskip}
\tema{Principios para evaluar el bienestar con discriminación}
\question ¿Cuáles son los principios básicos para evaluar cambios en el bienestar dada la discriminación?
\needspace{5\baselineskip}
\begin{solution}
Si la cantidad de equilibrio cae, entonces el bienestar total disminuye. Dadas dos demandas, puede que en equilibrio monopólico haya una ineficiencia por no servir cierto mercado; al discriminar y dejar de servir a más personas, se acentúa la ineficiencia. En caso de que la discriminación abra un mercado, por la misma lógica anterior debería aumentar el bienestar social.

\end{solution}

%-------------------------------------------------------------------- Demostración bienestar
\needspace{8\baselineskip}
\tema{Cotas al cambio en bienestar por discriminación de tercer grado}
\question Demuestre que el cambio en el bienestar al pasar de un monopolio no discriminante a un monopolio discriminante satisface
\[
  \sum_{i}(p_i - c)\Delta q_i \;\leq\; \Delta W \;\leq\; \sum_{i}(\bar p - c)\Delta q_i ,
\]
donde $\bar p$ es el precio uniforme, $p_i$ el precio discriminado en el mercado $i$, y $\Delta q_i = q_i(p_i) - q_i(\bar p)$ es el cambio en las ventas del mercado $i$ respecto a la situación sin discriminación. Interprete la desigualdad.
\needspace{5\baselineskip}
\begin{solution}
Primero establecemos la relación entre el precio, la demanda y el bienestar del consumidor. Con un precio $p$ se obtiene un excedente del consumidor $S_i(p) = \int_p^{p^*}q_i(s)\,ds$, el área bajo la curva de demanda desde $p$ hasta el precio donde la demanda es cero. Por lo tanto $S_i'(p) = -q_i(p)$ y $S_i''(p) = -q_i'(p) > 0$: el excedente del consumidor es convexo en el precio, lo que implica
\[
  S_i(p_i) - S_i(\bar p) \geq S_i'(\bar p)(p_i - \bar p) = -q_i(\bar p)(p_i - \bar p)
\]
\[
  \text{y} \qquad S_i(p_i) - S_i(\bar p) \leq S_i'(p_i)(p_i - \bar p) = -q_i(p_i)(p_i - \bar p).
\]
El cambio en el bienestar es
\[
  \Delta W = \sum_i\bigl[S_i(p_i) - S_i(\bar p)\bigr] + \sum_i\bigl[(p_i - c)q_i(p_i) - (\bar p - c)q_i(\bar p)\bigr].
\]
\emph{Cota inferior.} Usando la primera desigualdad, escribiendo $\bar q_i = q_i(\bar p)$ y $q_i = q_i(p_i)$:
\[
  \Delta W \geq \sum_i\bigl[-\bar q_i(p_i - \bar p) + (p_i - c)q_i - (\bar p - c)\bar q_i\bigr] = \sum_i\bigl[(p_i - c)q_i - (p_i - c)\bar q_i\bigr] = \sum_i(p_i - c)\Delta q_i .
\]
\emph{Cota superior.} Usando la segunda desigualdad:
\[
  \Delta W \leq \sum_i\bigl[-q_i(p_i - \bar p) + (p_i - c)q_i - (\bar p - c)\bar q_i\bigr] = \sum_i\bigl[(\bar p - c)q_i - (\bar p - c)\bar q_i\bigr] = \sum_i(\bar p - c)\Delta q_i .
\]
\textbf{Interpretación.} Se llega a los mismos principios discutidos antes: el cambio en bienestar está acotado por el cambio en las ventas valorado al margen de precio sobre costo. Si la discriminación lleva a un aumento de ventas por abrir un mercado ($\Delta q_i > 0$ en los mercados donde $p_i > c$), habrá un aumento en el bienestar; una condición necesaria para que el bienestar aumente es que la producción total aumente ($\sum\Delta q_i > 0$).

\end{solution}

%-------------------------------------------------------------------- Editorial
\needspace{8\baselineskip}
\tema{Discriminación intertemporal: tapa dura y tapa blanda}
\question Suponga una editorial con poder monopólico que está por lanzar un nuevo libro al mercado y que enfrenta dos tipos de compradores (hay una misma proporción de cada uno de los dos tipos). El primer tipo, que denominamos tipo 1, son pacientes y no les importa esperar hasta que salga la versión con tapa blanda. Su disposición máxima a pagar por el libro es $V_1 = 30$ con independencia de cuándo lo compran. El otro tipo de comprador, denominados tipo 2, son impacientes y quieren leer el libro lo antes posible. La disposición a pagar máxima de este grupo es $V_2 = 50e^{-\delta t}$, donde $t$ es el momento en el tiempo en que compra y lee el libro y $\delta$ es su tasa de descuento intertemporal. Así, estos compradores obtienen una utilidad máxima de 50 si compran y leen el libro inmediatamente ($t=0$). Cada consumidor compra solo un libro.

La editorial decide lanzar inmediatamente ($t=0$) un libro de tapa dura por un precio de $P_2$ para los lectores impacientes. Luego, tiene que decidir en qué momento del tiempo ($t$) lanzar una versión en tapa blanda a un precio menor $P_1$. El costo de imprimir un libro (sea de tapa dura o blanda) es irrelevante en comparación al precio, así que se asume que el costo marginal de producción es 0. Sin embargo, la editorial tiene una tasa de descuento para sus flujos igual a $r$ por cada unidad de tiempo $t$.
\begin{parts}
\part Escriba la restricción de racionalidad individual de los compradores tipo 1 (IR1), bajo el supuesto de que estos compran el libro de tapa blanda a un precio de $P_1$.
\needspace{5\baselineskip}
\begin{solution}
La disposición máxima de pago de los compradores tipo 1 es $V_1 = 30$, por lo que el precio del libro de tapa blanda no puede ser superior a esta cantidad:
\[
  IR_1:\quad V_1 - P_1 = 30 - P_1 \geq 0 .
\]

\end{solution}

\part Escriba la restricción de racionalidad individual de los compradores tipo 2 (IR2), bajo el supuesto de que estos compran el libro de tapa dura a un precio de $P_2$ (recordando que este libro se vende en $t=0$).
\needspace{5\baselineskip}
\begin{solution}
La disposición máxima de pago de los compradores tipo 2 en $t=0$ es $V_2|_{t=0} = 50$. Por lo tanto, el precio del libro de tapa dura no puede ser superior a 50:
\[
  IR_2:\quad V_2|_{t=0} - P_2 = 50 - P_2 \geq 0 .
\]

\end{solution}

\part Escriba la condición para que el lector tipo 1 prefiera esperar hasta que salga el libro de tapa blanda en vez de comprar el de tapa dura (IC1). Determine además la condición entre los precios para que se cumpla.
\needspace{5\baselineskip}
\begin{solution}
Para que los lectores tipo 1 prefieran comprar la tapa blanda en lugar de la tapa dura, el excedente neto al comprar el libro de tapa blanda debe ser mayor que el de tapa dura:
\[
  IC_1:\quad V_1 - P_1 \geq V_1 - P_2, \quad\text{es decir,}\quad 30 - P_1 \geq 30 - P_2 .
\]
Basta con que $P_2 \geq P_1$ para que se cumpla IC1.

\end{solution}

\part Escriba la condición para que el lector tipo 2 prefiera comprar inmediatamente el libro de tapa dura en vez de esperar hasta $t$ para comprar el libro de tapa blanda (IC2), teniendo en cuenta que su disposición a pagar cae con el tiempo.
\needspace{5\baselineskip}
\begin{solution}
Para que el lector tipo 2 prefiera comprar el libro de tapa dura inmediatamente, el excedente neto al comprarlo debe ser mayor que el de esperar por la versión en tapa blanda:
\[
  IC_2:\quad V_2|_{t=0} - P_2 \geq V_2|_{t} - P_1, \quad\text{es decir,}\quad 50 - P_2 \geq 50e^{-\delta t} - P_1 .
\]

\end{solution}

\part Suponga ahora que solo son restricciones activas (se cumplen con igualdad) IR1 e IC2. Escriba entonces las ganancias de la editorial y reemplace las restricciones para expresarlas solo en función de $t$ y los parámetros del problema.
\needspace{5\baselineskip}
\begin{solution}
Las ganancias de la editorial, considerando que los ingresos por la venta de los libros de tapa blanda ocurren en el momento $t$ y que los costos marginales de producción son 0, son $\pi = P_1e^{-rt} + P_2$. Reemplazando IR1 ($P_1 = 30$) e IC2 ($P_2 = 80 - 50e^{-\delta t}$):
\[
  \pi = 30e^{-rt} + 80 - 50e^{-\delta t}.
\]

\end{solution}

\part Obtenga la condición de primer orden para el $t$ óptimo y comente su respuesta en el contexto de discriminación de segundo grado.
\needspace{5\baselineskip}
\begin{solution}
Para maximizar las ganancias, derivamos $\pi$ respecto de $t$ e igualamos a cero:
\[
\begin{gathered}
  \frac{\partial\pi}{\partial t} = -30r\,e^{-rt} + 50\delta\,e^{-\delta t} = 0 \\
  \Longrightarrow\; t^* = \frac{\ln\!\bigl(\frac{3r}{5\delta}\bigr)}{r - \delta}.
\end{gathered}
\]
Este resultado indica que la editorial opta por retrasar el lanzamiento del libro de tapa blanda para capturar más valor de los lectores impacientes, lo cual es un ejemplo de discriminación de segundo grado en precios, donde se ofrece el mismo producto en diferentes versiones y tiempos para extraer mayor excedente del consumidor.

\end{solution}
\end{parts}

%-------------------------------------------------------------------- Comente: competencia y discriminación
\needspace{8\baselineskip}
\tema{Comente: discriminar bajo competencia}
\question \textbf{Comente.} En cualquier ambiente de competencia, las empresas van a preferir competir discriminando precios que aplicando un precio uniforme.
\needspace{5\baselineskip}
\begin{solution}
Falso. En monopolio siempre será preferible discriminar a fijar un precio uniforme. En el caso de competencia imperfecta, depende de la asimetría entre las firmas respecto de cuál es el mercado donde cobrarían un mayor precio si pudieran discriminar (como en el caso de AeroWeak, que prefiere que AeroStrong \emph{no} discrimine).

\end{solution}

%-------------------------------------------------------------------- Parque de diversiones
\needspace{8\baselineskip}
\tema{Tarifa en dos partes y tarjeta prepago: parque de diversiones}
\question Un parque de diversiones puede diferenciar dos tipos de clientes: niños y adultos. Los niños tienen una curva de demanda $Q_1 = 1000 - P$ y los adultos tienen una curva de demanda $Q_2 = 500 - \frac{P}{2}$. El costo marginal del parque es \$20, igual al costo medio. Se supone que la demanda por niños y adultos está conformada por 100 niños y 100 adultos.
\begin{parts}
\part Aplicar una tarifa en dos partes, cobrando una entrada al parque diferente a cada tipo de cliente y una tarifa igual para cada juego a que entre. Determine la utilidad obtenida por el monopolio.
\needspace{5\baselineskip}
\begin{solution}
La tarifa en dos partes consiste en fijar un precio por unidad igual al costo marginal (en este caso, \$20 por juego) y luego cobrar una tarifa de entrada que capture todo el excedente del consumidor. A $P = 20$, las cantidades demandadas por cada tipo son $Q_1 = 980$ (niños) y $Q_2 = 490$ (adultos). Los excedentes son
\[
  EC_{\text{niños}} = \tfrac{1}{2}\cdot 980\cdot 980 = 480\,200, \qquad EC_{\text{adultos}} = \tfrac{1}{2}\cdot 490\cdot 980 = 240\,100 .
\]
Por lo tanto, las tarifas de entrada son iguales al excedente de cada tipo (por persona): $\frac{480\,200}{100} = 4802$ para niños y $\frac{240\,100}{100} = 2401$ para adultos. El precio por juego es igual al costo marginal, por lo que toda la ganancia proviene de las tarifas de entrada. La utilidad total del monopolio es
\[
  \pi = 480\,200 + 240\,100 = 720\,300 .
\]

\end{solution}

\part El parque está analizando implementar una tarjeta de prepago con derecho a entrar a una cantidad de juegos diferente por tipo de cliente. ¿Cuánto debería cobrar por la tarjeta para cada tipo de cliente y cuántos juegos debería permitir?
\needspace{5\baselineskip}
\begin{solution}
Se calcula la disposición a pagar total de cada grupo con $P = 20$ (excedente más el gasto en juegos): niños, $480\,200 + 20\cdot 980 = 499\,800$; adultos, $240\,100 + 20\cdot 490 = 249\,900$. Por persona, la tarjeta cuesta $4998$ para niños (por $9{,}8$ juegos) y $2499$ para adultos (por $4{,}9$ juegos). La utilidad total es
\[
  \pi = 499\,800 + 249\,900 - 20\cdot 1470 = 720\,300 ,
\]
la misma que con la tarifa en dos partes.

\end{solution}

\part Compare las utilidades de los puntos anteriores con un monopolio que no discrimina tarifas (no cobra entrada, pero cobra una tarifa única por cada juego). ¿Cuál es la utilidad en ese caso?
\needspace{5\baselineskip}
\begin{solution}
La demanda agregada es $Q = Q_1 + Q_2 = 1500 - \frac{3}{2}P$, es decir, $P = 1000 - \frac{2}{3}Q$. El monopolio maximiza $\pi = (1000 - \tfrac{2}{3}Q)Q - 20Q$, con CPO $1000 - \tfrac{4}{3}Q - 20 = 0$, de donde $Q = 735$ y $P = 1000 - \tfrac{2}{3}\cdot 735 = 510$. Entonces $Q_{\text{niños}} = 490$, $Q_{\text{adultos}} = 245$ y la utilidad total es
\[
  \pi = 510\cdot 735 - 20\cdot 735 = 360\,150 ,
\]
bastante menor que los $720\,300$ con discriminación.

\end{solution}
\end{parts}

%-------------------------------------------------------------------- Discriminación perfecta tres tipos
\needspace{8\baselineskip}
\tema{Discriminación perfecta vs precio uniforme (tres tipos de consumidores)}
\question Considere un monopolista que cuenta con $N$ consumidores divididos en 3 tipos según sus demandas inversas: Tipo 1: $P = 10 - q$; Tipo 2: $P = 6 - 2q$; Tipo 3: $P = 3 - q$. El costo marginal del monopolista es constante e igual a 1. No hay costos fijos y el monopolista sabe que el porcentaje de consumidores del tipo 1 es $\lambda$ y del tipo 2 es $(1/2 - \lambda)$, con $0 < \lambda < 1/2$.
\begin{parts}
\part Encuentre las tarifas óptimas y los beneficios si el monopolista puede discriminar perfectamente a los distintos grupos de consumidores.
\needspace{5\baselineskip}
\begin{solution}
Con discriminación perfecta, el monopolista puede cobrar a cada tipo su excedente del consumidor como tarifa fija más el costo marginal por unidad ($p = c = 1$). Las cantidades son $q_1 = 9$, $q_2 = 2{,}5$, $q_3 = 2$, y los excedentes (tarifas fijas):
\[
  T_1 = \tfrac{1}{2}\cdot 9\cdot 9 = 40{,}5, \qquad T_2 = \tfrac{1}{2}\cdot 2{,}5\cdot 5 = 6{,}25, \qquad T_3 = \tfrac{1}{2}\cdot 2\cdot 2 = 2 .
\]
La proporción de consumidores tipo 3 es $1 - \lambda - (1/2 - \lambda) = 1/2$. Los beneficios son
\[
  \pi = N\bigl[40{,}5\lambda + 6{,}25(1/2 - \lambda) + 2\cdot\tfrac{1}{2}\bigr] = N\bigl[34{,}25\lambda + 4{,}125\bigr].
\]

\end{solution}

\part ¿Qué sucede con los beneficios cuando aumenta o disminuye $\lambda$? ¿Por qué se debe este resultado?
\needspace{5\baselineskip}
\begin{solution}
$\frac{\partial\pi}{\partial\lambda} = 34{,}25N > 0$: cuando $\lambda$ aumenta, los beneficios también aumentan. Esto se debe a que al aumentar $\lambda$ hay más consumidores tipo 1 (que tienen un mayor excedente) y menos tipo 2 (que generan menos excedente).

\end{solution}

\part Suponga que el gobierno decide prohibir este tipo de discriminación y obliga al monopolista a cobrar un precio único y uniforme para todos los consumidores. ¿Cuál es el precio que cobrará el monopolista? \emph{(Hint: suponga de aquí en adelante que $\lambda = 1/4$.)}
\needspace{5\baselineskip}
\begin{solution}
Con $\lambda = 1/4$ las proporciones son $\frac14$, $\frac14$ y $\frac12$. Las demandas individuales son $q_1 = 10 - p$, $q_2 = \frac{6-p}{2}$, $q_3 = 3 - p$. Si vende a los tres tipos ($p \leq 3$), la demanda agregada es
\[
  Q(p) = N\Bigl[\tfrac{1}{4}(10 - p) + \tfrac{1}{4}\cdot\tfrac{6-p}{2} + \tfrac{1}{2}(3 - p)\Bigr] = N\,\frac{38 - 7p}{8}.
\]
La CPO de $\max (p - 1)\frac{38 - 7p}{8}$ es $38 - 7p - 7(p - 1) = 0$, es decir $p = \frac{45}{14} \approx 3{,}21$. Como $p > 3$, los consumidores tipo 3 no compran. Si vende solo a los tipos 1 y 2 ($3 < p < 6$): $Q(p) = N\,\frac{26 - 3p}{8}$, y la CPO de $\max (p - 1)\frac{26 - 3p}{8}$ es $26 - 3p - 3(p - 1) = 0$, de donde
\[
  p = \frac{29}{6} \approx 4{,}83, \qquad \pi = N\Bigl(\tfrac{29}{6} - 1\Bigr)\frac{26 - 14{,}5}{8} = \frac{529}{96}N \approx 5{,}51N .
\]
Por lo tanto, solo compran los consumidores tipo 1 y tipo 2, a un precio de $4{,}83$. (Los beneficios son mucho menores que los $12{,}7N$ de la discriminación perfecta.)

\end{solution}

\part ¿Resulta la prohibición de discriminación una buena política óptima desde el punto de vista social?
\needspace{5\baselineskip}
\begin{solution}
La política del gobierno no es una buena política porque deja fuera del mercado a los individuos tipo 3. La discriminación perfecta permite que todos los tipos con disposición a pagar mayor al costo marginal puedan consumir, maximizando así el bienestar social. Sin discriminación, se genera una pérdida irrecuperable de eficiencia al excluir a los consumidores tipo 3 del mercado (y además los consumidores de los tipos 1 y 2 pagan un precio mayor al costo marginal).

\end{solution}
\end{parts}

%-------------------------------------------------------------------- Comente: 3er grado puede disminuir beneficios
\needspace{8\baselineskip}
\tema{Comente: discriminar en tercer grado vs precio lineal}
\question \textbf{Comente.} La discriminación de precios por indicadores (o de tercer grado) puede disminuir los beneficios del monopolista si lo comparamos con precios lineales. Verdadero o falso. Justifique.
\needspace{5\baselineskip}
\begin{solution}
Falso. El monopolista discriminador siempre puede fijar el mismo precio para cada grupo de consumidores y obtener los beneficios de los precios lineales del monopolista clásico. Por lo tanto, discriminar nunca puede dar menores beneficios que el precio uniforme.

\end{solution}

%-------------------------------------------------------------------- Tercer grado: 300-4p, 240-5p
\needspace{8\baselineskip}
\tema{Discriminación de tercer grado y precio uniforme impuesto por el gobierno}
\question Una empresa monopolista vende su producto para dos mercados distintos. Las demandas en esos mercados son $Q_1 = 300 - 4P$ y $Q_2 = 240 - 5P$. La función de costos de la empresa es $C(Q) = 10Q$.
\begin{parts}
\part Suponga que el monopolista puede practicar discriminación de tercer grado (es decir, cobrar precios distintos $P_1$ y $P_2$). Calcule los precios óptimos en cada uno de los mercados, las cantidades vendidas y el beneficio total del monopolista.
\needspace{5\baselineskip}
\begin{solution}
Se resuelve $\max\ (P_1 - 10)(300 - 4P_1) + (P_2 - 10)(240 - 5P_2)$. Las CPO son $300 - 4P_1 - 4(P_1 - 10) = 0 \Rightarrow P_1 = 42{,}5$ y $240 - 5P_2 - 5(P_2 - 10) = 0 \Rightarrow P_2 = 29$. Las cantidades vendidas son $Q_1 = 300 - 4\cdot 42{,}5 = 130$ y $Q_2 = 240 - 5\cdot 29 = 95$. Los beneficios son
\[
  \pi_1 = (42{,}5 - 10)\cdot 130 = 4225, \qquad \pi_2 = (29 - 10)\cdot 95 = 1805, \qquad \pi = 6030 .
\]

\end{solution}

\part Con el objetivo de disminuir la pérdida irrecuperable de eficiencia, promoviendo un mayor bienestar social, el gobierno decide obligar al monopolista a cobrar un precio único para los dos mercados. Calcule el precio óptimo, las cantidades vendidas y el beneficio total del monopolista bajo las nuevas reglas.
\needspace{5\baselineskip}
\begin{solution}
La demanda agregada es $Q = 540 - 9P$, y se resuelve $\max\ (P - 10)(540 - 9P)$: $540 - 9P - 9(P - 10) = 0 \Rightarrow P = 35$. Las cantidades son $Q_1 = 300 - 4\cdot 35 = 160$ y $Q_2 = 240 - 5\cdot 35 = 65$ (total $225$), y el beneficio total es $\pi = (35 - 10)\cdot 225 = 5625$.

\end{solution}

\part ¿Quién gana y quién pierde con la medida del gobierno?
\needspace{5\baselineskip}
\begin{solution}
Con la medida del gobierno, el bienestar del consumidor del mercado 1 aumenta (compra más, de 130 a 160, a menor precio, de 42,5 a 35): su excedente pasa de $2112{,}5$ a $3200$. El bienestar del consumidor del mercado 2 disminuye (compra menos, de 95 a 65, a mayor precio, de 29 a 35): su excedente pasa de $902{,}5$ a $422{,}5$. El monopolista pierde (sus beneficios caen de $6030$ a $5625$). Notar que la producción total es la misma en ambos casos ($225$), pero el bienestar total sube de $9045$ a $9247{,}5$, pues con un precio uniforme el bien se asigna mejor entre los consumidores.

\end{solution}
\end{parts}

%-------------------------------------------------------------------- Movistel
\needspace{8\baselineskip}
\tema{Discriminación de segundo grado: planes de internet (Movistel)}
\question Considere el siguiente problema. Hay dos tipos de consumidores de internet: los que juegan en línea y les importa la velocidad, y los que usan internet de manera general y les importa menos la velocidad. Suponga que hay 1.000 consumidores de cada tipo. La compañía proveedora de internet Movistel (un monopolista) ha estudiado la heterogeneidad de los consumidores y quiere ofrecer dos opciones de velocidad de bajada: un plan de 10 megas por segundo y un plan de 100 megas por segundo. Obviamente, todos los consumidores prefieren mayor velocidad pero, dada su heterogeneidad, están dispuestos a pagar cantidades distintas por cada opción. Los que juegan en línea están dispuestos a pagar \$10 por el plan de 100 megas y \$1 por el plan de 10 megas. Los que usan internet de manera general están dispuestos a pagar \$5 por el plan de 100 megas y \$0,75 por el plan de 10 megas. Suponga que la empresa no observa la disponibilidad a pagar de cada consumidor (aunque conoce las propiedades de la demanda). Suponga que el costo marginal es igual a 0.
\begin{parts}
\part Suponga que Movistel ofrece un único plan. ¿Qué plan ofrecerá y a qué precio? Calcule los beneficios de la empresa.
\needspace{5\baselineskip}
\begin{solution}
Hay que probar las 4 posibilidades:
\begin{enumerate}
  \item Plan de 100 megas, compra solo el tipo 1: precio 10 y beneficio $10\,000$.
  \item Plan de 100 megas, compran los tipos 1 y 2: precio 5 y beneficio $10\,000$.
  \item Plan de 10 megas, compra solo el tipo 1: precio 1 y beneficio $1000$.
  \item Plan de 10 megas, compran los tipos 1 y 2: precio $0{,}75$ y beneficio $1500$.
\end{enumerate}
De modo que el monopolista es indiferente entre ofrecer el plan de 100 megas a un precio de 10 o de 5, ya que en ambos casos obtiene un beneficio de $10\,000$.

\end{solution}

\part Suponga que Movistel ofrece un menú de precios (es decir, distintos planes). Escriba las condiciones de incentivos y participación para ambos grupos de consumidores. Calcule los precios óptimos y los beneficios.
\needspace{5\baselineskip}
\begin{solution}
Las condiciones de participación (PC) e incentivos (IC) para ambos grupos vienen dadas por:
\begin{align*}
  PC_1:&\quad 10 - p_{100} \geq 0, & IC_1:&\quad 10 - p_{100} \geq 1 - p_{10}, \\
  PC_2:&\quad 0{,}75 - p_{10} \geq 0, & IC_2:&\quad 0{,}75 - p_{10} \geq 5 - p_{100}.
\end{align*}
Los precios óptimos son $p_{10} = 0{,}75$ (se extrae todo el excedente del tipo 2) y $p_{100} = 9{,}75$ (el máximo que hace que el tipo 1 no prefiera el plan de 10 megas: $10 - p_{100} = 1 - 0{,}75$). El beneficio obtenido es $1000\cdot 9{,}75 + 1000\cdot 0{,}75 = 10\,500$.

\end{solution}

\part Determine si Movistel ofrece el menú de precios.
\needspace{5\baselineskip}
\begin{solution}
Sí, porque el beneficio obtenido ($10\,500$) es mayor que sin diferenciación ($10\,000$).

\end{solution}
\end{parts}
